\documentclass[11pt]{article}
% Shared print layout. The web's provisional design and MathJax are independent.
\usepackage[utf8]{inputenc}
\usepackage[T1]{fontenc}
\usepackage[spanish,es-nodecimaldot]{babel}
\usepackage{lmodern}
\usepackage{amsmath,amssymb}
\usepackage{graphicx}
\usepackage{geometry}
\geometry{a4paper,left=20mm,right=20mm,top=25mm,bottom=22mm,headheight=16pt,headsep=8mm}
\usepackage{fancyhdr,lastpage}
\usepackage{enumitem}
\usepackage{xcolor}
\usepackage{pgfplots}
\pgfplotsset{compat=1.18}
\usepgfplotslibrary{fillbetween}
\usepackage[hidelinks,pdfencoding=auto]{hyperref}
\hypersetup{pdflang=es-ES,pdfauthor={},pdfsubject={Matemáticas Aplicadas a las Ciencias Sociales II}}
\pagestyle{fancy}
\fancyhf{}
\fancyhead[L]{\small Matemáticas Aplicadas a las Ciencias Sociales II}
\fancyhead[R]{\small 2.\textsuperscript{o} Bachillerato}
\fancyfoot[L]{\footnotesize Subida de nota T1 | Edición revisada 1.0}
\fancyfoot[R]{\footnotesize Página \thepage\ de \pageref{LastPage}}
\renewcommand{\headrulewidth}{0.4pt}
\setlength{\parindent}{0pt}
\setlength{\parskip}{0.5em}
\setlist{itemsep=0.3em,topsep=0.3em}
\newcommand{\apartado}[1]{\subsection*{#1}}
\pgfplotsset{
  revision/.style={
    width=0.94\linewidth,height=7.5cm,
    axis lines=middle,grid=major,grid style={gray!18},
    tick label style={font=\small},label style={font=\small},
    xlabel={$x$},ylabel={$y$},samples=160,
    legend style={at={(0.5,-0.16)},anchor=north,draw=none,font=\small},
    clip=true
  }
}

\hypersetup{pdftitle={Subida de nota T1 - Soluciones revisadas}}
\begin{document}
\section*{Subida de nota T1: soluciones}
\section*{Ejercicio 1(a): continuidad y derivabilidad}
El enunciado del examen utiliza una \textbf{rama racional}, no una rama lineal:
\[f(x)=\begin{cases}\dfrac{a}{x-1},&x<0,\\x^2-bx-1,&x\geq0.\end{cases}\]
Para que sea derivable en $0$, primero debe ser continua:
\[f(0)=-1,\qquad\lim_{x\to0^-}f(x)=-a,\qquad\lim_{x\to0^+}f(x)=-1.\]
Así, $-a=-1$ y obtenemos $a=1$.

En el interior de cada intervalo,
\[f'(x)=\begin{cases}-\dfrac{a}{(x-1)^2},&x<0,\\2x-b,&x>0.\end{cases}\]
Bajo la condición de continuidad, las derivadas laterales en la unión son
\[f'_-(0)=-a=-1,\qquad f'_+(0)=-b.\]
La igualdad exige $-1=-b$, es decir,
\[\boxed{a=1,\qquad b=1}.\]
\textbf{Comprobación por definición.} Para estos parámetros y un incremento $t$,
\[
\frac{f(t)-f(0)}{t}=\begin{cases}\dfrac{1}{t-1},&t<0,\\t-1,&t>0.\end{cases}
\]
Ambos cocientes tienden a $-1$; por tanto, $f'(0)=-1$.

\textbf{Aclaración editorial.} El par $a=1$, $b=-1$ del archivo de soluciones original corresponde a otra función, cuya rama izquierda es $a(x-1)$. No resuelve el enunciado de este examen.
\clearpage
\section*{Ejercicio 1(b): área en $[0,1]$}
Este apartado fija \textbf{otros parámetros}: $a=1$, $b=2$. En todo $[0,1]$ se utiliza la segunda rama:
\[f(x)=x^2-2x-1=(x-1)^2-2.\]
Como $0\leq(x-1)^2\leq1$ en ese intervalo, se tiene $-2\leq f(x)\leq-1$. La gráfica está por debajo del eje de abscisas, de modo que
\begin{align*}
A&=\int_0^1|f(x)|\,dx=-\int_0^1(x^2-2x-1)\,dx\\
&=\left[-\frac{x^3}{3}+x^2+x\right]_0^1
=-\frac13+1+1=\boxed{\frac53\text{ unidades}^2}.
\end{align*}
\textbf{Comprobación.} La altura del recinto está entre 1 y 2 y su base mide 1; por tanto $1\leq A\leq2$, coherente con $5/3$.
\begin{center}
\begin{tikzpicture}
\begin{axis}[revision,xmin=-0.2,xmax=1.3,ymin=-2.4,ymax=0.5]
\addplot[name path=curve,thick,domain=0:1] {x^2-2*x-1};
\addplot[name path=zero,draw=none,domain=0:1] {0};
\addplot[gray!25] fill between[of=curve and zero];
\addplot[thick,domain=0:1] {x^2-2*x-1};
\addplot[dashed] coordinates {(1,0) (1,-2)};
\end{axis}
\end{tikzpicture}
\end{center}
La figura representa exclusivamente el recinto solicitado; la rama para $x<0$ no interviene en esta área.
\clearpage
\section*{Ejercicio 2(a): estudio y representación}
\[g(x)=\frac{x}{x+3}=1-\frac{3}{x+3},\qquad D_g=\mathbb R\setminus\{-3\}.\]
\textbf{Continuidad y asíntotas.} Es continua en cada intervalo de su dominio. En $x=-3$ hay una discontinuidad infinita:
\[\lim_{x\to-3^-}g(x)=+\infty,\qquad\lim_{x\to-3^+}g(x)=-\infty.\]
La asíntota vertical es $x=-3$. Como $\lim_{x\to\pm\infty}g(x)=1$, la horizontal es $y=1$.

\textbf{Monotonía.} La derivada es
\[g'(x)=\frac{(x+3)-x}{(x+3)^2}=\frac{3}{(x+3)^2}>0.\]
Crece estrictamente en $(-\infty,-3)$ y en $(-3,+\infty)$; no hay extremos relativos. No es creciente en todo el dominio considerado conjuntamente: $-4<-2$, pero $g(-4)=4>-2=g(-2)$.

\textbf{Cortes.} Numerador nulo en $0$, con denominador distinto de cero. Ambos ejes se cortan en $(0,0)$.

\textbf{Ampliación conservada del solucionario.} $g''(x)=-6/(x+3)^3$: curvatura hacia arriba para $x<-3$ y hacia abajo para $x>-3$. No hay punto de inflexión en $-3$, pues no pertenece al dominio. La curvatura no es un requisito adicional del enunciado.
\begin{center}
\begin{tikzpicture}
\begin{axis}[revision,height=6.2cm,xmin=-8,xmax=5,ymin=-4,ymax=5]
\addplot[thick,domain=-8:-3.06] {x/(x+3)};
\addplot[thick,domain=-2.94:5] {x/(x+3)};
\addplot[dashed] coordinates {(-3,-4) (-3,5)};
\addplot[dotted,domain=-8:5] {1};
\addplot[only marks,mark=*,mark size=1.7pt] coordinates {(0,0)};
\end{axis}
\end{tikzpicture}
\end{center}
Las ramas se dibujan por separado, sin evaluar el denominador nulo.
\clearpage
\section*{Ejercicio 2(b): recta tangente}
El punto de tangencia y la pendiente son
\[P=(2,g(2))=\left(2,\frac25\right),\qquad m=g'(2)=\frac3{25}.\]
Aplicamos la ecuación punto-pendiente:
\[y-\frac25=\frac3{25}(x-2).\]
Al ordenar,
\[\boxed{y=\frac3{25}x+\frac4{25}}.\]
\textbf{Comprobación.} En $x=2$, la recta da $(6+4)/25=2/5$; su pendiente coincide con $g'(2)$.
\begin{center}
\begin{tikzpicture}
\begin{axis}[revision,xmin=-0.5,xmax=5,ymin=-0.3,ymax=0.9]
\addplot[thick,domain=-0.5:5] {x/(x+3)};
\addlegendentry{$g(x)=x/(x+3)$}
\addplot[dashed,domain=-0.5:5] {3*x/25+4/25};
\addlegendentry{Tangente: $y=3x/25+4/25$}
\addplot[only marks,mark=*,mark size=2pt] coordinates {(2,0.4)};
\end{axis}
\end{tikzpicture}
\end{center}
\clearpage
\section*{Ejercicio 3: primitiva y dominio}
\[h(x)=x^2+e^{2x}-\frac1{1-x},\qquad x\neq1.\]
Integramos por términos. Para el racional, el cambio $u=1-x$, $du=-dx$, da
\[\int\frac1{1-x}\,dx=-\ln|1-x|+C.\]
Como en $h$ este término lleva signo negativo,
\begin{align*}
\int h(x)\,dx
&=\frac{x^3}{3}+\frac{e^{2x}}2-\bigl(-\ln|1-x|\bigr)+C\\
&=\boxed{\frac{x^3}{3}+\frac{e^{2x}}2+\ln|1-x|+C}.
\end{align*}
\textbf{Dominio y constantes.} La fórmula vale en cada intervalo que no cruce $1$. En $(-\infty,1)$ y $(1,+\infty)$ pueden elegirse constantes independientes $C_-$ y $C_+$. Si solo se pide una primitiva en el dominio, se pueden tomar ambas iguales a cero.

\textbf{Comprobación por derivación:}
\[
\frac{d}{dx}\left(\frac{x^3}{3}+\frac{e^{2x}}2+\ln|1-x|\right)
=x^2+e^{2x}-\frac1{1-x}=h(x).
\]
El valor absoluto es necesario para escribir una expresión real válida también cuando $x>1$.

\textbf{Corrección editorial.} El paso $\int1/(1-x)\,dx=x$ del original es incorrecto: la derivada de $x$ es $1$. El resultado correcto es $-\ln|1-x|+C$.
\end{document}
